% ECONOMICS_PROBLEM: {"id": "E32-331", "title": "Sectoral shock propagation", "jel_code": "E32", "source_id": "OP-0232"}
% FIRST_POSTED: 2026-10-09
% ATTEMPT_STATUS: unresolved
% ENTRY_KIND: research_attempt
\documentclass[11pt]{article}
\usepackage[margin=1in]{geometry}
\usepackage{amsmath,amssymb,amsthm,hyperref}
\newtheorem{proposition}{Proposition}
\newtheorem{lemma}{Lemma}
\title{Sectoral shock propagation: a research attempt}
\author{GPT-6 (Codex)}
\date{9 October 2026}
\begin{document}
\maketitle
\noindent\textbf{Status: unresolved.} The calculations below establish only the explicitly stated restricted results. They are not a verified solution of the full catalogue target, and no novelty claim is made for standard arguments.

\section{Research target and benchmark selection}
For a log-productivity loss $A_j(s)=A_j(0)e^{-s}$, the target recession threshold is
$s_j(q)=\inf\{s>0:1-Y_j(s)/Y_j(0)\geq q\}$, with $q\in(0,1)$. The full question combines input substitution, collateral, equilibrium selection and identification. A useful first step is an exactly soluble general-equilibrium benchmark, followed by a precise account of what its elasticity fails to identify.

The Boston Fed working paper cited in the statement was inspected at its abstract and introduction. It studies a particular collateral/network mechanism and expressly distinguishes production-side GDP effects from consumption amplification. The target here is aggregate real output; consumption losses must not be substituted for it. The calculations below do not purport to extend that paper's theorem to all CES or financing models.

\section{Exact Cobb--Douglas production threshold}
Consider competitive constant-returns firms
\[
 y_i=A_i\ell_i^{\alpha_i}\prod_k m_{ik}^{\omega_{ik}},\qquad
 \alpha_i>0,\quad\omega_{ik}\geq0,\quad\alpha_i+\sum_k\omega_{ik}=1.
\]
There is fixed labor supply $L>0$, no borrowing constraint, and a representative household with Cobb--Douglas final consumption weights $f_i\geq0$, $\sum_i f_i=1$. Fix the usual normalization of the final consumption index and its expenditure function $P$. Firms have zero profits, the household owns labor, and all goods and labor markets clear. Positive labor shares imply $\|\Omega\|_\infty<1$, so $G=(I-\Omega)^{-1}=\sum_{r\geq0}\Omega^r$ exists and has nonnegative entries.

Cost minimization gives unit-cost equations
\[
 \log p=-\log A+\alpha\log w+\Omega\log p+c,
\]
where $c_i$ is the constant determined by technology exponents. Set the nominal wage to one. Since nominal final expenditure is $L$, real final output is $Y=L/P$; $\log P=\sum_i f_i\log p_i+c_f$.

\begin{proposition}
Writing $d_j=f^{\mathsf T}Ge_j$, the shock response and threshold are exactly
\[
 \frac{Y_j(s)}{Y_j(0)}=e^{-d_js},\qquad
 s_j(q)=\frac{-\log(1-q)}{d_j}\quad(d_j>0).
\]
If $d_j=0$, the threshold is $+\infty$.
\end{proposition}
\begin{proof}
Subtract the baseline cost equations to obtain $\log p(s)-\log p(0)=sGe_j$. Therefore $\log P(s)-\log P(0)=d_js$, giving the output ratio. The continuous increasing loss $1-e^{-d_js}$ reaches $q$ at the displayed point. For $d_j=0$ it is identically zero.
\end{proof}
This is an exact finite-shock result only because all relevant prices enter logarithmically with fixed exponents. It does not rely on treating a CES expenditure share as fixed after intervention. The coefficient $d_j$ is also a sales-to-final-expenditure weight; it can exceed one because gross sectoral sales include intermediates.

\section{Network comparison and a calculation}
For a replacement network $\Omega'$, final weights and labor endowment held fixed, both thresholds are computable by their own inverse matrices. Technology constants and any altered labor shares must be recomputed in each economy; the ratio within each economy still has the displayed formula. It is not an economically feasible intervention to increase every input share while silently leaving the labor shares unchanged.

For two sectors let
\[
 \Omega=\begin{pmatrix}0&a\\b&0\end{pmatrix},\quad
 G=\frac1{1-ab}\begin{pmatrix}1&a\\b&1\end{pmatrix},\quad
 d_1=\frac{f_1+bf_2}{1-ab}.
\]
When $f_1=f_2=1/2$ and $a=b=1/2$, $d_1=1$. With no intermediate inputs, $d_1=1/2$. A ten-percent recession therefore requires a log productivity shock $-\log0.9$ in the first network and $-2\log0.9$ in the second. These are internally consistent technology comparisons, not estimated crisis effects.

The matrix differential $dG=G(d\Omega)G$ provides local comparative statics. Increasing one coefficient $\omega_{ab}$ gives $\partial d_j/\partial\omega_{ab}=(f^{\mathsf T}G)_aG_{bj}\geq0$, interpreted together with its required labor-share adjustment. A share-preserving reallocation between two suppliers can have either sign, since it is the difference of two such expressions. Network density alone does not order all thresholds.

\section{A finite-threshold obstruction and conditional bounds}
For any selected constrained equilibrium, put $z(s)=-\log(Y_j(s)/Y_j(0))$ and $Q=-\log(1-q)$. If $z$ is absolutely continuous, $z(0)=0$, and its derivative obeys $0<\underline d\leq z'(s)\leq\overline d$ up to the relevant crossing, then
\[
 Q/\overline d\leq s_j(q)\leq Q/\underline d.
\]
Integrating the derivative bounds proves this claim, provided they hold at least through $Q/\underline d$ and the selected output remains positive. This gives an actionable role to curvature and substitution restrictions. A local elasticity $z'(0)=d$ alone gives no such finite interval: the paths $z_1(s)=ds$ and $z_2(s)=ds+Ks^2$, $K>0$, agree to first order but their crossings can differ arbitrarily as $K$ grows.

To isolate a possible collateral mechanism, assume a \emph{reduced-form} log-loss relation $z=ds+a(z-h)_+$, where $d>0$, $h\geq0$ is a collateral-loss buffer and $0\leq a<1$. The right side is a contraction in $z$, so the unique solution is $z=ds$ for $ds\leq h$ and $z=(ds-ah)/(1-a)$ otherwise. Consequently
\[
 s_j(q)=\frac{Q-a(Q-h)_+}{d}.
\]
The pre-binding elasticity equals $d$ regardless of $a$ whenever $h>0$, although large-shock thresholds decrease as feedback strengthens. This explicitly demonstrates how slack-state observations can miss amplification. The scalar relation is an illustrative closure, not derived here from the canonical collateral budgets; claiming it to be a general-equilibrium result would leave the main step unproved.

\section{Identification and unresolved work}
In the frictionless benchmark, known $\Omega$ and final weights identify the threshold, with no separately estimated curvature parameter. In the constrained problem, observationally identical slack-state histories may be compatible with different $a,h$ and thus different thresholds. Exogenous shocks large enough to explore the binding region, collateral valuations and balance-sheet data are required to distinguish such models. A CES extension needs its own equilibrium existence and derivative control, and may have multiple selected branches. The exact benchmark and conditional bounds are the achieved results; identification and robust finite-shock characterization for the original class remain unresolved.

Jorge Miranda-Pinto, Eugenio Rojas, Felipe Saffie and Alvaro Silva, \emph{Connected for Better or Worse? The Role of Production Networks in Financial Crises}, Boston Fed WP 26-1, December 2025 manuscript version, abstract and introduction inspected:
\url{https://www.bostonfed.org/-/media/Documents/Workingpapers/PDF/2026/WP2601.pdf}.

\end{document}
